% Part: first-order-logic
% Chapter: sequent-calculus
% Section: structural-rules

\documentclass[../../../include/open-logic-section]{subfiles}

\begin{document}

\iftag{FOL}
      {\olfileid{fol}{seq}{srl}}
      {\olfileid{pl}{seq}{srl}}

\olsection{قواعد ساختاری}

همچنین به چند قاعده نیاز داریم که امکان بازآرایی !!{sentence}s را در
سمت چپ و راست دنباله فراهم کنند. چون قواعد منطقی ایجاب می‌کنند که در
مقدمه، !!{sentence}sیی که قاعده بر آن‌ها عمل می‌کند یا در منتهی‌الیه
چپ باشند یا در منتهی‌الیه راست، به قاعدهٔ «تبادل» نیاز داریم تا
!!{sentence}s را به جایگاه درست منتقل کنیم. گاهی نیز مهم است که بتوان
از میان !!{sentence}s، دو مورد یکسان را در یکی ادغام کرد و
!!a{sentence} را به هر
یک از دو سو افزود.

\subsection{تضعیف}

\begin{defish}
\Axiom$ \Gamma \fCenter \Delta $
\RightLabel{\LeftR{\Weakening}}
\UnaryInf$ !A, \Gamma \fCenter \Delta$
\DisplayProof
\hfill
\Axiom$ \Gamma \fCenter \Delta$
\RightLabel{\RightR{\Weakening}}
\UnaryInf$ \Gamma \fCenter \Delta, !A$
\DisplayProof
\end{defish}

\subsection{انقباض}

\begin{defish}
\Axiom$ !A, !A, \Gamma \fCenter \Delta $
\RightLabel{\LeftR{\Contraction}}
\UnaryInf$ !A, \Gamma \fCenter \Delta$
\DisplayProof
\hfill
\Axiom$ \Gamma \fCenter \Delta, !A, !A$
\RightLabel{\RightR{\Contraction}}
\UnaryInf$ \Gamma \fCenter \Delta, !A$
\DisplayProof
\end{defish}

\subsection{تبادل}

\begin{defish}
\Axiom$ \Gamma, !A, !B, \Pi \fCenter \Delta $
\RightLabel{\LeftR{\Exchange}}
\UnaryInf$ \Gamma, !B, !A, \Pi \fCenter \Delta$
\DisplayProof
\hfill
\Axiom$ \Gamma \fCenter \Delta, !A, !B, \Lambda$
\RightLabel{\RightR{\Exchange}}
\UnaryInf$ \Gamma \fCenter \Delta, !B, !A, \Lambda$
\DisplayProof
\end{defish}

زنجیره‌ای از استنتاج‌های تضعیف، انقباض و تبادل را اغلب با خط‌های
استنتاج دوتایی نشان می‌دهیم.

قاعدهٔ زیر که «برش» نام دارد، به معنای دقیق کلمه ضروری نیست، اما
استفادهٔ دوباره و ترکیب !!{derivation}s را بسیار آسان‌تر می‌کند.
\begin{defish}
\[
\Axiom$ \Gamma \fCenter \Delta, !A$
\Axiom$ !A, \Pi \fCenter \Lambda $
\RightLabel{\Cut}
\BinaryInf$ \Gamma, \Pi \fCenter \Delta, \Lambda$
\DisplayProof
\]
\end{defish}

\end{document}
